% Small-budget regime (C2 sweep). Compressed again for plan v3 (2026-09-07): two paragraphs to one, full
% numbers in the artifact. Numbers unchanged. v2 in sections/backup_v2_2026-09-07/.
\subsection{The small-budget regime}\label{sec:sweep}

To observe the mechanism where it acts, we reran the six prompt classes (900 prompts, three trajectories each, $T_{\max}=200$, temperature 1) at $k \in \{0.1, 0.15, 0.25, 0.5, 1\}$ with both baselines on the same prompts and seeds. With raw prompts the constraint is active on $68$ to $74\%$ of decode steps for $k \le 0.25$, $49\%$ at $k=0.5$, and $5.5\%$ at $k=1$; no generation at $k \le 0.25$ is identical to the risky model's and $4\%$ are at $k=1$, against $57\%$ between $k=3$ and $k=5$. The anchor writes the opening of every answer, $39$ tokens at $k=0.1$, $27$ at the implementation's default $k=0.15$, and $4$ at $k=1$, matching $\lfloor \delta_{\mathrm{init}}/k \rfloor$: at the budgets where the certificate is informative, a fifth of a $200$-token answer is written by a 1.8B model that has not read the book, and Section~\ref{sec:utility} measures what that costs. The event-level tails are heavier here too: the realised log-likelihood ratio exceeds $K$ for $13\%$ of attack-split trajectories at $k=0.1$, $33\%$ at $k=0.5$, and $10\%$ at $k=1$, with 95\% anytime-valid intervals excluding zero at every budget, so the budgets at which the certificate is informative are the budgets at which up to a third of outputs individually exceed it. Copying metrics are flat across every budget and both baselines (ROUGE-L $0.058$ to $0.062$, near-verbatim recall $0$), because this risky model does not reproduce the passages; on this workload a sweep of $k$ measures utility loss, not protection. The chat template changes the size but not the shape, and none of the $37{,}800$ trajectories violated an invariant (\texttt{regime\_table.csv} in the artifact).
